\subsection{Iterated derivatives}

1.7.1. Let \(f\) be a function defined in a neighborhood of a point \(x_{0}\) of E, taking values in F. If \(f\) is differentiable in a neighborhood of \(x_{0}\) , its derivative Df is a mapping from a neighborhood of \(x_{0}\) in the polynormed space \(\mathcal{L}(\mathrm{E};\mathrm{F})\) of continuous linear mappings from E into F. Let \(p\) be an integer \(\geqslant 2\) :

we say that \(f\) is \(p\) times differentiable at \(x_0\) if \(f\) is differentiable in a neighborhood of \(x_0\) and if its derivative \(Df\) is \((p - 1)\) times differentiable at \(x_0\) . We then define the \(p\)-th derivative of \(f\) at \(x_0\); this is the continuous \(p\)-linear map \(D^p f(x_0)\) from \(E^p\) into \(F\) , defined by:

\[
\mathrm{D} ^ {p} f (x _ {0}). (h _ {1}, \dots , h _ {p}) = (\mathrm{D} (\mathrm{D} ^ {p - 1} f) (x _ {0}). h _ {1}). (h _ {2}, \dots , h _ {p}).
\]

We also set \(D^{0}f = f\) and \(D^{1}f = Df\) . If \(f\) is \(p\) times differentiable at \(x_{0}\) and if \(q\) and \(s\) are two integers such that \(q + s = p\) , with \(s > 0\), then \(f\) is \(q\) times differentiable in a neighborhood of \(x_{0}\) , the function \(D^{q}f\) (with values in \(\mathcal{L}_{q}(E; F)\) ) is \(s\) times differentiable at \(x_{0}\) , and we have:

\[
\mathrm{D} ^ {q + s} f (x _ {0}). (h _ {1}, \dots , h _ {q + s}) = (\mathrm{D} ^ {s} (\mathrm{D} ^ {q} f) (x _ {0}). (h _ {1}, \dots , h _ {s})) \cdot (h _ {s + 1}, \dots , h _ {q + s})
\]

a relation which, by abuse of notation, is written in the form:

\[
\mathbf {D} ^ {q + s} f = \mathbf {D} ^ {s} \mathbf {D} ^ {q} f.
\]

1.7.2. Let \((\mathbf{E}_{i})_{1\leqslant i\leqslant n}\) be closed vector subspaces of E, such that E is the topological direct sum of the \(E_{i}\) . One then defines, when it exists, the iterated partial derivative \(D_{i_{1}}\ldots D_{i_{m}}f\) of a map f from a neighborhood of \(x_{0}\in E\) into F; it is a continuous multilinear map from

\[
\mathbf {E} _ {i _ {1}} \times \dots \times \mathbf {E} _ {i _ {m}}
\]

into F, defined by induction on the integer \(m \geqslant 1\) as follows: if \(D_{i_{2}} \ldots D_{i_{m}} f(x)\) exists in a neighborhood of \(x_{0}\) and has a partial derivative with respect to \(E_{i_{1}}\) , then \(D_{i_{1}} \ldots D_{i_{m}} f(x_{0})\) is given by:

\[
\mathrm{D} _ {i _ {1}} \dots \mathrm{D} _ {i _ {m}} f (x _ {0}). (h _ {1}, \dots , h _ {m}) = (\mathrm{D} _ {i _ {1}} (\mathrm{D} _ {i _ {2}} \dots \mathrm{D} _ {i _ {m}} f) (x _ {0}). h _ {1}). (h _ {2}, \dots , h _ {m})
\]

for \(h_{k} \in E_{i_{k}}\) .

If \(f\) is \(m\) times differentiable at \(x_{0}\) , then the partial derivative \(\mathbf{D}_{i_{1}}\ldots\mathbf{D}_{i_{m}}f(x_{0})\) exists and is equal to the restriction of \(\mathbf{D}^{m}f(x_{0})\) to the subspace \(\mathbf{E}_{i_{1}}\times\cdots\times\mathbf{E}_{i_{m}}\) of \(\mathbf{E}^{m}\) . Consequently, \(\mathbf{D}^{m}f(x_{0})\) is completely determined by the iterated partial derivatives of order \(m\) at \(x_{0}\) .

1.7.3. Suppose that E is finite-dimensional and let \((e_{1},\ldots,e_{n})\) be a basis of E. Set \(E_{i}=Ke_{i}\) and let f be a map from a neighborhood of \(x_{0}\) into F. If the partial derivative \(D_{i_{1}}\ldots D_{i_{m}}f(x_{0})\) (with the notations of 1.7.2) exists, one sets:

\[
\partial_ {i _ {1}} \dots \partial_ {i _ {m}} f (x _ {0}) = \mathbf {D} _ {i _ {1}} \dots \mathbf {D} _ {i _ {m}} f (x _ {0}). (e _ {i _ {1}}, \dots , e _ {i _ {m}}).
\]

